\documentclass[11pt]{article}
% Shared print layout. The web's provisional design and MathJax are independent.
\usepackage[utf8]{inputenc}
\usepackage[T1]{fontenc}
\usepackage[spanish,es-nodecimaldot]{babel}
\usepackage{lmodern}
\usepackage{amsmath,amssymb}
\usepackage{graphicx}
\usepackage{geometry}
\geometry{a4paper,left=20mm,right=20mm,top=25mm,bottom=22mm,headheight=16pt,headsep=8mm}
\usepackage{fancyhdr,lastpage}
\usepackage{enumitem}
\usepackage{xcolor}
\usepackage{pgfplots}
\pgfplotsset{compat=1.18}
\usepgfplotslibrary{fillbetween}
\usepackage[hidelinks,pdfencoding=auto]{hyperref}
\hypersetup{pdflang=es-ES,pdfauthor={},pdfsubject={Matemáticas Aplicadas a las Ciencias Sociales II}}
\pagestyle{fancy}
\fancyhf{}
\fancyhead[L]{\small Matemáticas Aplicadas a las Ciencias Sociales II}
\fancyhead[R]{\small 2.\textsuperscript{o} Bachillerato}
\fancyfoot[L]{\footnotesize Sistemas de ecuaciones II | Edición revisada 1.0}
\fancyfoot[R]{\footnotesize Página \thepage\ de \pageref{LastPage}}
\renewcommand{\headrulewidth}{0.4pt}
\setlength{\parindent}{0pt}
\setlength{\parskip}{0.5em}
\setlist{itemsep=0.3em,topsep=0.3em}
\newcommand{\apartado}[1]{\subsection*{#1}}
\pgfplotsset{
  revision/.style={
    width=0.94\linewidth,height=7.5cm,
    axis lines=middle,grid=major,grid style={gray!18},
    tick label style={font=\small},label style={font=\small},
    xlabel={$x$},ylabel={$y$},samples=160,
    legend style={at={(0.5,-0.16)},anchor=north,draw=none,font=\small},
    clip=true
  }
}

\hypersetup{pdftitle={Sistemas de ecuaciones II - Soluciones revisadas}}
\begin{document}
\section*{Sistemas de ecuaciones II: soluciones}
\section*{Ejercicio 1: sustitución regresiva}
La ecuación matricial equivale a
\[
\begin{cases}2x+y-z=4,\\3y+5z=7,\\z=1.\end{cases}
\]
Resolvemos de abajo hacia arriba.
\subsection*{1. Obtener $z$}
La tercera ecuación da directamente $z=1$.
\subsection*{2. Obtener $y$}
\[3y+5\cdot1=7\quad\Longrightarrow\quad3y=2\quad\Longrightarrow\quad y=\frac23.\]
\subsection*{3. Obtener $x$}
\[2x+\frac23-1=4\quad\Longrightarrow\quad2x=4+\frac13=\frac{13}3.\]
Así,
\[\boxed{(x,y,z)=\left(\frac{13}6,\frac23,1\right)}.\]
\textbf{Comprobación en las tres ecuaciones:}
\begin{align*}
2\cdot\frac{13}6+\frac23-1&=\frac{13}3+\frac23-1=4,\\
3\cdot\frac23+5\cdot1&=2+5=7,\\
z&=1.
\end{align*}
La solución es única: el determinante de la matriz triangular es $2\cdot3\cdot1=6\neq0$.
\clearpage
\section*{Ejercicio 2: empresa de logística}
Sean $S$, $M$ y $L$ las cantidades de envíos de paquetes pequeños, medianos y grandes, respectivamente. Deben ser enteros no negativos.

La diferencia del enunciado es \textbf{medianos menos grandes}, de modo que
\[
\begin{cases}S+M+L=150,\\10S+15M+20L=2\,450,\\S=2(M-L).\end{cases}
\]
\subsection*{1. Relacionar $L$ con $M$}
Sustituyendo $S=2(M-L)$ en el total,
\[2(M-L)+M+L=150\quad\Longrightarrow\quad3M-L=150.\]
Por tanto, $L=3M-150$.
\subsection*{2. Obtener $M$ con los ingresos}
\begin{align*}
10\cdot2(M-L)+15M+20L&=2\,450,\\20M-20L+15M+20L&=2\,450,\\35M&=2\,450.
\end{align*}
Así, $M=70$.
\subsection*{3. Obtener los otros dos tipos}
\[L=3\cdot70-150=60,\qquad S=2(70-60)=20.\]
\[\boxed{S=20,\qquad M=70,\qquad L=60}.\]
\textbf{Comprobación completa:}
\begin{align*}
20+70+60&=150,\\10\cdot20+15\cdot70+20\cdot60&=200+1\,050+1\,200=2\,450,\\2(70-60)&=20=S.
\end{align*}
Las cantidades son enteras no negativas y $M\geq L$, coherente con la relación dada.
\clearpage
\section*{Ejercicio 3(a): discusión según $\lambda$}
\[
A=\begin{pmatrix}1&1&1\\2&3&\lambda\\4&5&7\end{pmatrix},\qquad
b=\begin{pmatrix}3\\6\\10\end{pmatrix}.
\]
Calculamos el determinante:
\begin{align*}
\det A&=(21-5\lambda)-(14-4\lambda)+(10-12)\\
&=5-\lambda.
\end{align*}
\subsection*{Caso 1: $\lambda\neq5$}
El determinante no es cero. Por tanto,
\[\operatorname{rg}(A)=\operatorname{rg}(A|b)=3.\]
El sistema es \textbf{compatible determinado}: tiene solución única.
\subsection*{Caso 2: $\lambda=5$}
Para decidir la compatibilidad, transformamos la matriz ampliada mediante $F_2\leftarrow F_2-2F_1$ y $F_3\leftarrow F_3-4F_1$:
\[
\left(\begin{array}{ccc|c}1&1&1&3\\2&3&5&6\\4&5&7&10\end{array}\right)
\ \longrightarrow\ 
\left(\begin{array}{ccc|c}1&1&1&3\\0&1&3&0\\0&1&3&-2\end{array}\right).
\]
Ahora $F_3\leftarrow F_3-F_2$ da
\[
\left(\begin{array}{ccc|c}1&1&1&3\\0&1&3&0\\0&0&0&-2\end{array}\right).
\]
La última fila exige $0=-2$, lo que es imposible. Los rangos son
\[\operatorname{rg}(A)=2<3=\operatorname{rg}(A|b).\]
El sistema es \textbf{incompatible}: no tiene solución.
\[
\boxed{\begin{array}{ll}\lambda\neq5:&\text{solución única (SCD)},\\\lambda=5:&\text{sin solución (SI)}.\end{array}}
\]
No hay valores del parámetro que produzcan infinitas soluciones.
\clearpage
\section*{Ejercicio 3(b): resolución para $\lambda=0$}
Como $0\neq5$, existe una única solución. El sistema es
\[
\begin{cases}x+y+z=3,\\2x+3y=6,\\4x+5y+7z=10.\end{cases}
\]
Aplicamos las mismas eliminaciones que antes:
\[
\begin{cases}y-2z=0,\\y+3z=-2.\end{cases}
\]
Restando la primera a la segunda,
\[5z=-2\quad\Longrightarrow\quad z=-\frac25.\]
Sustituyendo,
\[y=2z=-\frac45,\qquad x=3-y-z=3+\frac45+\frac25=\frac{21}5.\]
Por tanto,
\[\boxed{(x,y,z)=\left(\frac{21}5,-\frac45,-\frac25\right)}.\]
\textbf{Comprobación en las ecuaciones originales con $\lambda=0$:}
\begin{align*}
\frac{21}5-\frac45-\frac25&=3,\\
2\cdot\frac{21}5+3\left(-\frac45\right)&=\frac{42-12}5=6,\\
4\cdot\frac{21}5+5\left(-\frac45\right)+7\left(-\frac25\right)&=\frac{84-20-14}5=10.
\end{align*}
Las tres ecuaciones se cumplen. Estos valores negativos son admisibles en el sistema algebraico; no representan cantidades de envíos.
\end{document}
